\chapter{Informes de avance y cambios}

Este anexo resume la evolución del proyecto a lo largo de sus principales fases de trabajo. Los informes recogen el avance funcional, las decisiones tomadas y los cambios relevantes respecto a la planificación inicial.

\section{Informe de avance 1: definición del alcance y análisis inicial}

\subsection{Trabajo realizado}

Durante la fase inicial se definió el problema que motiva el proyecto: la dispersión de la información sobre eventos y actividades en Sevilla. Se analizó el contexto de uso, se estudiaron soluciones existentes y se delimitó el alcance de \textit{Sevillaneando} como aplicación móvil participativa orientada al descubrimiento de eventos.

También se identificaron los actores principales del sistema, los objetivos del proyecto y los objetivos funcionales de la aplicación.

\subsection{Cambios y decisiones}

Se decidió enfocar el producto en eventos locales y emergentes, diferenciándolo de directorios generales mediante funcionalidades sociales, eventos privados, recomendaciones y rutas de eventos.

\section{Informe de avance 2: requisitos, planificación y diseño conceptual}

\subsection{Trabajo realizado}

En esta fase se elaboró el documento de requisitos, incluyendo requisitos de información, requisitos funcionales, requisitos no funcionales y reglas de negocio. Además, se preparó la planificación del proyecto mediante fases, sprints, estimación de esfuerzo, EDT y diagrama temporal.

Se diseñaron los primeros mockups y el modelo conceptual del sistema, sentando la base para el diseño técnico posterior.

\subsection{Cambios y decisiones}

Se adoptó una metodología ágil adaptada a un proyecto unipersonal, con seguimiento mediante backlog, sprints y registro de horas. La estimación se ajustó a una dedicación aproximada de 300 horas, conforme a la carga esperada del TFG.

\section{Informe de avance 3: arquitectura e implementación base}

\subsection{Trabajo realizado}

Se definió la arquitectura cliente-servidor del sistema, separando la aplicación móvil, el backend y la base de datos. Se implementó la estructura inicial del backend con NestJS, la aplicación móvil con React Native y Expo, y la persistencia con PostgreSQL y PostGIS.

Durante esta fase se desarrollaron las funcionalidades básicas de autenticación, gestión de usuarios, eventos, categorías y consulta de datos.

\subsection{Cambios y decisiones}

Se consolidó el uso de TypeScript en todo el sistema para mantener coherencia entre frontend y backend. También se decidió utilizar servicios externos como Firebase Authentication y Cloudinary para reducir complejidad propia en autenticación y almacenamiento de imágenes.

\section{Informe de avance 4: funcionalidades sociales, mapas y recomendaciones}

\subsection{Trabajo realizado}

Se incorporaron funcionalidades avanzadas: mapa de eventos, geolocalización, filtrado por distancia, chat de evento, mensajería privada, módulo de amistad basado en seguimiento, valoraciones, recomendaciones personalizadas y rutas de eventos.

El sistema empezó a comportarse como una plataforma social completa, no solo como un catálogo de eventos.

\subsection{Cambios y decisiones}

El desarrollo del sistema de rutas y recomendaciones requirió ajustar el alcance de algunas funcionalidades secundarias. Se priorizaron las funcionalidades centrales de descubrimiento, participación y personalización frente a extras inicialmente previstos.

\section{Informe de avance 5: scraping, moderación y despliegue}

\subsection{Trabajo realizado}

Se integró el scraping de eventos desde fuentes externas, incluyendo procesamiento de datos, geocodificación y control de duplicados. También se desarrolló el flujo de moderación de eventos públicos, con estados de aprobación y rechazo, y se preparó el despliegue del backend y la distribución de la aplicación móvil.

Además, se documentó el proceso mediante manuales de desarrollo, despliegue y usuario.

\subsection{Cambios y decisiones}

La complejidad del scraping fue superior a la prevista debido a la variabilidad de las fuentes externas, la extracción de direcciones, el tratamiento de fechas y los límites de peticiones. Como consecuencia, se descartó la implementación final de funcionalidades opcionales como la generación de carteles con IA y la moderación automática de imágenes mediante NSFWJS.

\section{Informe de avance 6: pruebas, validación y cierre}

\subsection{Trabajo realizado}

En la fase final se ejecutaron pruebas unitarias, pruebas funcionales y pruebas de aceptación sobre los flujos principales de la aplicación. También se recogieron observaciones de usuarios piloto y se corrigieron incidencias detectadas durante la validación.

Se completó la memoria, se prepararon los anexos y se revisó la coherencia documental del trabajo.

\subsection{Cambios y decisiones}

Las pruebas con usuarios piloto permitieron detectar incidencias específicas en Android, especialmente relacionadas con el mapa y la disposición visual de algunos elementos. Estas correcciones introdujeron una desviación de esfuerzo, pero mejoraron la estabilidad de la aplicación en un entorno multiplataforma.

\section{Resumen de cambios relevantes}

\begin{itemize}
    \item Incorporación de CI/CD con GitHub Actions para automatizar comprobaciones y despliegues.
    \item Ajuste del alcance por la complejidad del scraping y del sistema de rutas.
    \item Descarte justificado de funcionalidades opcionales no prioritarias.
    \item Refuerzo de pruebas en Android tras detectar incidencias en dispositivos reales.
    \item Ampliación de la documentación final con manuales, glosario, actas e informes de avance y cambios.
\end{itemize}
